\section{Method}\label{sec:method}
For each position $p$ and each angular bin (binned in sky tangents in world coordinates) we define the opacity
\begin{equation}
\begin{aligned}
\lambda_p&=-\ln\!\left[\frac{n_p}{(t_p/t_{\mathrm{sky}})\,n_{\mathrm{sky}}}\right],\\
\sigma_{\lambda,p}^2&\simeq\frac{1}{n_p}+\frac{1}{n_{\mathrm{sky}}}.
\end{aligned}
\label{eq:opacity}
\end{equation}
where $n$ are track counts and $t$ live times. The open-sky ratio normalizes the detector response, assuming it remained stable between runs. The live-time ratio fixes the zero of $\lambda$, so no free offsets are fitted in the nominal voxel inversion. Nominal fit weights use binwise standard deviations estimated from \InputsOpacityWeightReplicas{} Poisson replicas; Eq.~\ref{eq:opacity} gives the analytic approximation. Bins without sky data are masked. Figure~\ref{fig:opacity} shows the opacity maps.

\begin{figure*}[t]
\centering
\includegraphics[width=0.88\textwidth]{opacity.pdf}
\caption{Opacity seen from the two room positions, and open-sky counts. The beam shadows appear as repeated stripes on a smooth background from the slab.}\label{fig:opacity}
\end{figure*}

For the three-dimensional image we write $\boldsymbol\lambda=A\boldsymbol q$, where $q_j$ is the opacity per unit length in voxel $j$ and $A_{ij}$ is the path length of bin $i$ through voxel $j$, averaged over the detector aperture. A track with local detector tangents $(t_x,t_y)$ can only pass through all layers within an area of size $(W-|t_x|D)\times(W-|t_y|D)$, with $W$ the active width and $D$ the layer gap; we sample this area with sub-rays and rotate it with the detector azimuth. The system is solved with non-negative iterative updates \citep{gilbert1972} and total-variation smoothing in the horizontal plane only \citep{rudin1992}. We do not smooth vertically, because the height information comes from parallax and smoothing would blur it. Voxels seen by few rays are damped to suppress edge artefacts. The voxel size is $\ReconstructionSpacing\,\mathrm{m}$; this is the grid, not the resolution, which is much coarser in $z$.

The pose of pos1 (floor offset, separation and azimuth) is fitted by requiring that a reconstruction from one position predicts the data of the other. To avoid locking onto the beam pitch, the fit is averaged over two grids shifted by half a voxel. It gives a floor offset $(\PoseX,\PoseY)\,\mathrm{m}$, a separation of $\PoseBaseline\,\mathrm{m}$ and an azimuth of $\PoseAz^\circ$. The absolute scale remains conditional on the rough floor-offset prior $\InputsFloorOffsetPrior\,\mathrm{m}$, search bounds and reconstruction objective. The finite detector footprint and voxel lattice also introduce length scales; the fitted pose is not an independent survey. On-site ceiling measurements were not used to choose it.
